% Part: first-order-logic
% Chapter: natural-deduction
% Section: quantifier-rules

\documentclass[../../../include/open-logic-section]{subfiles}

\begin{document}

\olfileid{fol}{ntd}{qrl}

\olsection{Regels voor kwantoren}

\subsection{Regels voor $\lforall$ (voor alle)}

\begin{defish}
\AxiomC{$!A(a)$}
\RightLabel{\Intro{\lforall}}
\UnaryInfC{$\lforall[x][\Atom{!A}{x}]$}
\DisplayProof
\hfill
\AxiomC{$\lforall[x][\Atom{!A}{x}]$}
\RightLabel{\Elim{\lforall}}
\UnaryInfC{$!A(t)$}
\DisplayProof
\end{defish}

In deze regels voor~$\lforall$ is $t$ een gesloten term, dus een term
zonder variabelen. De letter $a$ is !!a{constant} met een bijzondere
beperking: zij mag niet voorkomen in de conclusie
$\lforall[x][!A(x)]$ en ook niet in een aanname die !!{undischarged} is
in de !!{derivation} die eindigt met de premisse~$!A(a)$. We noemen
$a$ de \emph{eigenvariabele} van de afleidingsstap
\Intro{\lforall}.\footnote{We noemen $a$ hier om historische redenen
een ``eigenvariabele'', hoewel zij in de bovenstaande regel een
constante is.}

\subsection{Regels voor $\lexists$ (er bestaat)}

\begin{defish}
\AxiomC{$\Atom{!A}{t}$}
\RightLabel{\Intro{\lexists}}
\UnaryInfC{$\lexists[x][\Atom{!A}{x}]$}
\DisplayProof
\hfill
\AxiomC{$\lexists[x][\Atom{!A}{x}]$}
\AxiomC{[$\Atom{!A}{a}$]$^n$}
\DeduceC{$!C$}
\DischargeRule{\Elim{\lexists}}{n}
\BinaryInfC{$!C$}
\DisplayProof
\end{defish}

Ook hier is $t$ een gesloten term. De !!a{constant}~$a$ mag niet
voorkomen in de premisse $\lexists[x][!A(x)]$, in de conclusie~$!C$ of
in een aanname die !!{undischarged} is in de !!{derivation}s die naar
de twee premissen leiden. Alleen de aannamen $!A(a)$ vormen daarop een
uitzondering. We noemen $a$ de \emph{eigenvariabele} van de
afleidingsstap \Elim{\lexists}.

We kunnen beide beperkingen in één regel samenvatten. Een
eigenvariabele mag niet in de premissen voorkomen en ook niet in een
aanname die !!{undischarged} is in de !!{derivation}s die naar de
premissen van de afleidingsstap \Intro{\lforall} of \Elim{\lexists}
leiden. Dit heet de \emph{eigenvariabelevoorwaarde}.

\begin{explain}
We gebruiken de volgende notatie. Als $!A$ !!a{formula} is waarin de
!!{variable}~$x$ vrij voorkomt, schrijven we dat als~$!A(x)$. In die
context kort $!A(t)$ de substitutie~$\Subst{!A}{t}{x}$ af. Daardoor
kunnen we de regel $\Intro\lexists$ ook als volgt schrijven:
\begin{prooftree}
\AxiomC{$\Subst{!A}{t}{x}$}
\RightLabel{\Intro{\lexists}}
\UnaryInfC{$\lexists[x][!A]$}
\end{prooftree}
De term $t$ kan al in~$!A$ staan; $!A$ kan bijvoorbeeld
$\Atom{\Obj P}{t,x}$ zijn. Daarom is uit~$\Atom{\Obj P}{t,t}$ de
formule $\lexists[x][\Atom{\Obj P}{t,x}]$ afleiden een correcte
toepassing van~$\Intro\lexists$. Bij deze stap mag men één of meer
voorkomens van~$t$ in de premisse door de gebonden !!{variable}~$x$
``vervangen''; niet alle voorkomens hoeven te worden vervangen. Voor
$\Intro\lforall$ en~$\Elim\lexists$ is de beperking strenger: de
!!{constant}~$a$ mag niet in~$!A$ voorkomen. Daarom kan men met
$\Intro\lforall$ niet correct uit $\Atom{\Obj P}{a,a}$ de formule
$\lforall[x][\Atom{\Obj P}{a,x}]$ afleiden.
\end{explain}

\begin{explain}
Vergelijk nu de twee soorten regels. Bij \Intro{\lexists} en
\Elim{\lforall} mag de term~$t$ willekeurig zijn; daar hoeven we dus
geen extra voorwaarde te controleren. Bij \Elim{\lexists} en
\Intro{\lforall} geldt wel de eigenvariabelevoorwaarde: de
!!{constant}~$a$ mag nergens in de conclusie of in een
!!{undischarged} aanname voorkomen. Deze beperking houdt het systeem
correct. Dat betekent dat het alleen !!{sentence}s !!{derive} uit
!!{undischarged} aannamen waaruit zij volgen. Zonder de beperking zou
de volgende ongeldige stap zijn toegestaan:
\begin{prooftree}
  \AxiomC{$\lexists[x][!A(x)]$}
  \AxiomC{$\Discharge{!A(a)}{1}$}
  \RightLabel{*\Intro{\lforall}}
  \UnaryInfC{$\lforall[x][!A(x)]$}
  \RightLabel{\Elim{\lexists}}
  \BinaryInfC{$\lforall[x][!A(x)]$}
\end{prooftree}
Maar $\lexists[x][!A(x)] \Entails/ \lforall[x][!A(x)]$. Een
bestaansbewering rechtvaardigt dus niet zonder meer een algemene
bewering.

Ten slotte mogen we volgens de eliminatieregels voor kwantoren alleen
gesloten termen voor !!{variable}s substitueren. Daarom is elke
!!{formula} die uit een verzameling !!{sentence}s kan worden afgeleid,
zelf !!a{sentence}.
\end{explain}
\end{document}
